\documentclass[../comorg-part1.tex]{subfiles}
\begin{document}

\section{Instruction Organisation in Memory}
Instructions are stored in memory as machine codes (binary patterns).\newline
ARM instructions are human-readable representations of machine code.

\medskip
\noindent ARM instructions are 32 bits long.

\begin{figure}[H]
    \centering
    \includegraphics[width=0.7\textwidth]{images/instruction-format.png}
    \caption*{ARM Instruction format (32 bits long)}
\end{figure}
\begin{itemize}
    \item Opcode specifies the operation to be carred out (Chapter 5 - Instruction Set)
    \item Operands specify the data or location of data to be used / stored. (Chapter 4 - Addreessing Modes).
\end{itemize}

\subsection{ARM Programmer's Model}
Programmer's model is an abstract, simplified view of the processor.\newline
It includes just the parts a programmer needs to write and run code.
\begin{figure}[H]
    \centering
    \includegraphics[width=0.3\textwidth]{images/cpu-registers.png}
    \caption*{Visible registers in User Mode of the ARM processor}
\end{figure}
\begin{itemize}
    \item There are 16 32-bit registers.
    \item R13 (SP) stores address of some space in memory, that temporarily stores away register information that may be needed.\newline
    Used to maintain full descending stack etc.
    \item R14 (LP) Used as Link Register when calling of subroutines.\newline
    For example, when $\texttt{BL}$ (Branch with Link) called, R14 stores PC (R15) value so that the program can return when subroutine is completed.\newline
    R14 can also be a general purpose register at other times.
    \item R15 (PC) stores the start address of the next instruction to be \textbf{fetched}.\newline
    PC will automatically increment by the length of the instruction executed (4 for ARM CPU).\newline
    Jump or branch instruction will alter the sequential execution.\newline
    The value of PC is the address of the current instruction being executed plus 8 bytes (due to the fetch-decode-execute pipeline).
    \item CPSR stores the 4 condition code flags.\newline
    Bit 31: N (Negative), Bit 30: Z (Zero), Bit 29: C (Carry out), Bit 28: V (Overflow)
\end{itemize}

\subsection{Execution Cycle}
The CPU repeatedly performs the Fetch-Decode-Execute cycle.

\noindent Considering the instruction $\texttt{LDR R1, [R0]}$, and where the instruction and data are stored:
\begin{figure}[H]
    \centering
    \includegraphics[width=0.3\textwidth]{images/executioncycle-eg.png}
    % \caption*{}
\end{figure}
We get this execution cycle:
\begin{enumerate}
    \item Fetch cycle - Instruction
    \begin{figure}[H]
        \centering
        \includegraphics[width=0.8\textwidth]{images/fetch-cycle-instruction.png}
        \caption*{PC points to address of next instruction, the opcode fed into IR.}
    \end{figure}
    \item Fetch cycle - Operand
    \begin{figure}[H]
        \centering
        \includegraphics[width=0.8\textwidth]{images/fetch-cycle-operand.png}
        \caption*{Another fetch required to retrieve operand in CPU.}
    \end{figure}
    \item Execute Cycle - Load
    \begin{figure}[H]
        \centering
        \includegraphics[width=0.8\textwidth]{images/execute-cycle-load.png}
        \caption*{The operand fetched from memory is loaded into memory.}
    \end{figure}
\end{enumerate}

\noindent Takeaways:
\begin{enumerate}
    \item Different instructions require multiple accesses to memory, hence taking different number of clock cycles.
    \item \noindent von Neumann bottleneck - System performance limited by data traffic bandwidth between CPU and memory, because
data transfer on the external bus is slower than within CPU's internal bus. Hence keeping regularly used operands in the CPU registers may help reduce memory access.
    \item \noindent Keeping instructions and data in separate memories (i.e. Harvard architecture)
can make instructions execute in more regular cycles (using parallel fetches).
\end{enumerate}

\end{document}